\documentclass{article}
\usepackage[submission]{colm2026_conference}

\usepackage[utf8]{inputenc}
\usepackage[T1]{fontenc}
\usepackage{microtype}
\usepackage[hypertexnames=false]{hyperref}
\usepackage{url}
\usepackage{booktabs}
\usepackage{amsmath}
\usepackage{amssymb}
\usepackage{graphicx}
\usepackage{algorithm}
\usepackage{algpseudocode}
\usepackage{lineno}

\definecolor{darkblue}{rgb}{0, 0, 0.5}
\hypersetup{colorlinks=true, citecolor=darkblue, linkcolor=darkblue, urlcolor=darkblue}

\title{Minimal Execution Patches for Causal Debugging of Tool-Using Agents}

\author{Anonymous Authors}

\begin{document}

\ifcolmsubmission
\linenumbers
\fi

\maketitle

\begin{abstract}
Tool-using language agents often fail because one intermediate action corrupts
the downstream execution state, yet standard benchmark evaluation usually
reveals only the final reward and a raw trajectory trace. We propose execution
intervention for post-hoc debugging: a methodology that treats failed agent
trajectories as replayable programs. Given a failed run with checkpointed
environment state, the method restores an intermediate snapshot, applies a
small structured patch, replays execution, and checks whether the benchmark
outcome flips from failure to success. The smallest successful intervention is
useful both as a repair and as a causal autopsy of the root failure. We
instantiate this protocol on real tau-style environments with deterministic
state, saved benchmark trajectories, and benchmark-native evaluators. The
current system supports structured patch families, bounded search-cost
accounting, strict versus oracle continuation analysis, and explicit rerun
baselines including retry from scratch and retry from a localized snapshot. On
controlled synthetic failures, the prototype supports exact localization
metrics in addition to recovery. On imported natural failures, the current
artifact-backed evidence already includes strict non-oracle recoveries in both
retail and airline settings, including continuation-based recovery after an
identity-lookup repair and deletion-based recovery for harmful mutating
actions. The frozen workshop corpus now contains $80$ replay-failing natural
cases and $24$ controlled synthetic failures across retail and airline. On the
updated workshop artifacts, structured strict replay recovers $27/80$ natural
failures, the oracle upper bound reaches $35/80$, and the strongest simple retry
baseline recovers $15/80$; on the synthetic corpus, heuristic localization
recovers $24/24$ with perfect top-1 localization. These results motivate
replayable execution intervention as a practical debugging protocol for LLM
agents: more causal than prompt-only retry, more interpretable than
unconstrained regeneration, and naturally aligned with benchmark evaluation.
\end{abstract}

\section{Introduction}

Many agent benchmarks now measure multi-step interaction with tools, APIs,
databases, browsers, or code environments. When such an agent fails, the usual
options are limited: inspect the trace manually, rerun with a stronger model,
or prompt the model to reflect and retry. These approaches can improve task
success, but they do not cleanly answer a debugging question: which
intermediate decision actually caused the failure?

Our central claim is that replayable environments let us ask a stronger causal
question. If we can checkpoint and restore execution state at each step, then a
failed trajectory is not merely a transcript. It is a stateful program that
can be paused, forked, edited, and re-executed. This enables a simple but
useful search problem: find the smallest intervention that flips the final
benchmark outcome.

This submission is methodology-first. We do not claim benchmark-wide repair at
this stage. Instead, we contribute a concrete protocol built around
checkpointed replay, structured patch families, bounded search cost, and
autopsy generation, together with a now-frozen workshop corpus over saved
tau-style failures in retail and airline domains.

Existing work mostly predicts or explains where an agent failed from traces. We
instead search for the smallest executable intervention that causally flips the
benchmark outcome under replay, yielding both a repair and a root-cause
autopsy. That difference matters: our unit of analysis is not only a suspect
step, but a benchmark-native reward flip produced by checkpointed replay,
minimal patch size, explicit search-cost accounting, and strict non-oracle
continuation after the intervention.

\section{Method}

\begin{figure}[t]
\centering
\fbox{%
\parbox{0.96\linewidth}{%
\small
\textbf{Failed trace}
$\rightarrow$
\textbf{restore ranked snapshot}
$\rightarrow$
\textbf{propose structured patch}
$\rightarrow$
\textbf{replay under strict continuation}
$\rightarrow$
\textbf{benchmark evaluator}
$\rightarrow$
\textbf{minimal successful patch}
$\rightarrow$
\textbf{autopsy report}

\medskip
The key distinction is executable counterfactual repair: we intervene on saved
environment state, not only on textual explanations of the trace.
}%
}
\caption{Replayable execution intervention treats a failed agent trajectory as a
stateful program that can be paused, patched, replayed, and scored by the
original benchmark evaluator.}
\label{fig:execution_intervention_pipeline}
\end{figure}

For a failed trajectory, we log per-step prompt context, chosen action, tool
result, and pre/post execution snapshots. A patch candidate specifies a target
step, a patch family, and a small structured payload. Applying a patch means:
\begin{enumerate}
  \item restore the saved pre-step snapshot;
  \item edit the action or context according to the patch payload;
  \item resume execution under a bounded continuation policy;
  \item evaluate the replayed rollout with the original benchmark evaluator.
\end{enumerate}

\begin{algorithm}[t]
\caption{Minimal Patch Recovery via Checkpointed Replay}
\label{alg:minimal_patch_recovery}
\begin{algorithmic}[1]
\Require failed trajectory $\tau$, ranked steps $\mathcal{T}$, patch families $\mathcal{F}$, search budget $B$
\State $p^\star \gets \varnothing$
\For{each $t \in \mathcal{T}$ in ranked order}
  \State restore the saved pre-step snapshot at step $t$
  \For{$f \in \mathcal{F}$}
    \For{$p \in \textsc{ProposePatches}(\tau, t, f, B)$}
      \State replay from the restored snapshot with patch $p$
      \State continue under a bounded strict continuation policy
      \State evaluate the replayed rollout with the benchmark reward
      \If{$R(\tau'(p)) = 1$}
        \If{$p^\star = \varnothing$ or $p$ is smaller/cheaper than $p^\star$}
          \State $p^\star \gets p$
        \EndIf
      \EndIf
    \EndFor
  \EndFor
\EndFor
\State \Return winning patch $p^\star$ and its autopsy report, or \texttt{not recovered}
\end{algorithmic}
\end{algorithm}

The current system supports structured intervention families including tool
argument edits, tool-call replacement, insertion, deletion, context edits, and
continuation replacement. Search returns the smallest successful patch under a
budget over localization, proposal, and replay. The output is not only a
success flag: it is also an autopsy artifact with the root-cause step, patch
family, patch size, total search cost, and before/after state fragments.

The workshop implementation also includes three simple rerun baselines:
\begin{itemize}
  \item retry from scratch,
  \item retry from the localized snapshot without structured patching,
  \item raw continuation from the localized snapshot while preserving the
  original failing action.
\end{itemize}

\paragraph{Cost accounting and proxies.}
To evaluate search efficiency across deterministic and stochastic execution
paths, we use a dual cost-accounting protocol. For online baselines that call
OpenRouter, we report provider-returned token usage from billing metadata. For
the deterministic patch engine, we optimize an internal analytical cost model
$C(P)$ rather than a provider-specific tokenizer. This makes the search-cost
metric reproducible across runs and hardware settings, and avoids coupling the
core method to a changing external tokenizer.
\begin{equation}
C(P) = C_{\text{local}} + C_{\text{patch}} + C_{\text{replay}} + C_{\text{cont}} .
\end{equation}
Here $C_{\text{local}} = \rho \cdot \left|S_{\text{ranked}}\right|$ penalizes rollback and
step ranking depth, with default $\rho = 8$. $C_{\text{patch}}$ is a structural
penalty determined by patch family plus a length proxy from standardized text
serialization: for example, tool-argument edits carry a base penalty of $10$,
tool-call replacement $12$, context edits $6$, and deletions $4$, each plus a
whitespace-token length term where applicable. $C_{\text{replay}}$ and
$C_{\text{cont}}$ capture replay and continuation overhead, including bounded
continuation proposals generated during deterministic replay. We therefore
report efficiency as a unified search-cost accounting layer, while explicitly
distinguishing actual API tokens from analytical proxy units when discussing
specific methods.

\section{Workshop Evaluation Scope}

The workshop path stays tau-native and focuses on retail plus airline. The
frozen corpus currently contains $40$ retail plus $40$ airline natural
failures, together with $12$ retail plus $12$ airline synthetic failures. The
main headline setting is strict non-oracle replay: after patching a saved
failure, the system must continue from the replayed state without reusing
oracle future actions. Oracle continuation is retained only as an upper-bound
analysis.

The workshop comparison frame is methodology-first with early real-environment
evidence rather than benchmark-scale repair claims. The intended ablation set
includes no repair, random candidate / random localization, reverse search,
tool-deletion-only search, continuation-replacement-only search, and the
oracle suffix upper bound. The corpus scale has already moved from the earlier
thirty-two-case smoke setting to a frozen $80$-case natural corpus. The final
artifact-backed workshop tables now show $27/80$ strict natural recoveries
($33.8\%$), $35/80$ oracle-upper-bound recoveries ($43.8\%$), and retry-baseline
recoveries of $15/80$ from scratch, $7/80$ from a localized snapshot, and
$3/80$ for raw continuation from the failing action. On the frozen $24$-case
synthetic corpus, heuristic search reaches $24/24$ recovery with perfect
top-1 localization.

The workshop artifact pipeline supports building natural and synthetic corpora,
running structured patch search and retry baselines, exporting tables and
figure CSVs, and generating reviewer-facing bundles. The artifact layer uses
resumable JSONL outputs and sharded failure corpora so the frozen workshop
release can be regenerated without eagerly materializing a multi-gigabyte
corpus in memory. The artifact package includes saved oracle autopsies, paper
tables, figure CSVs, and a curated three-case study bundle over the frozen
artifacts.

\subsection*{Workshop Results}

Table~\ref{tab:workshop_natural_results} summarizes the main natural-failure
comparison. Under strict non-oracle replay, structured patch search recovers
$27/80$ failures, compared with $15/80$ for retry from scratch, $7/80$ for
retry from a localized snapshot, and $3/80$ for raw continuation from the
failing step. The oracle upper bound
recovers $35/80$, leaving an eight-case strict-versus-oracle gap that is more
about continuation quality than fault localization alone. The retry baselines
also highlight a practical context-window bottleneck: naive rerun-style
continuation consumes between $199{,}323$ and $349{,}895$ provider-reported
tokens on the frozen corpus, showing that unguided long-horizon retry quickly
becomes expensive and increasingly brittle as the execution history grows.

\begin{table}[t]
\centering
\small
\begin{tabular}{lrrrr}
\toprule
Method & Recov. & SRR & Cost metric$^{*}$ & Median$^{*}$ \\
\midrule
Structured patch search (strict) & 27/80 & 0.338 & 8751 & 61.0 \\
Oracle upper bound & 35/80 & 0.438 & 9698 & 61.0 \\
Retry from scratch & 15/80 & 0.188 & 199323 & 2561.0 \\
Retry from localized snapshot & 7/80 & 0.088 & 349895 & 3827.0 \\
Raw continuation from snapshot & 3/80 & 0.038 & 228922 & 3343.0 \\
\bottomrule
\end{tabular}
\caption{Workshop natural-failure results on the frozen $80$-case retail plus
airline corpus. Structured patch search outperforms the simple rerun-style
baselines under the same saved-failure protocol.}
\vspace{0.25em}

\begin{minipage}{0.96\linewidth}
\footnotesize
$^{*}$ Structured patch search and the oracle upper bound report the analytical
search-cost model described in the text. Retry baselines report actual
provider-returned API token usage. These values should therefore be interpreted
as an efficiency accounting layer rather than a single billing-equivalent
tokenizer.
\end{minipage}
\label{tab:workshop_natural_results}
\end{table}

The natural results also show a clear domain imbalance. In strict mode, the
current workshop system recovers $22/40$ airline failures but only $5/40$
retail failures. Under the oracle upper bound, airline rises only slightly to
$23/40$, while retail rises to $12/40$. This pattern is consistent with our
core claim: better continuation helps, but the benchmark-native replay
interventions are already exposing where simple deletion-style repair works and
where richer continuation or context repair is still needed.

An important empirical finding is that the strict natural recoveries in the
frozen workshop corpus are currently deletion-dominated rather than broadly
distributed across all patch families. We view this as informative rather than
embarrassing. Many replay-failing airline cases contain harmful extra or
mutating actions, so minimal deletion is exactly the kind of executable
counterfactual that should flip the evaluator outcome if the suspected action
is truly causal. The retail gap to the oracle upper bound suggests a different
failure mode: those cases appear less blocked by local harmful actions and more
blocked by downstream continuation quality, context selection, or multi-step
identity repair.

\begin{table}[t]
\centering
\small
\begin{tabular}{lrr}
\toprule
Successful natural patch family & Strict & Oracle \\
\midrule
Tool deletion & 27 & 33 \\
Tool call with oracle suffix & 0 & 2 \\
\bottomrule
\end{tabular}
\caption{Successful natural patch-family distribution in the frozen workshop
release. The strict leaderboard is deletion-dominated, suggesting that many
recoverable airline failures arise from harmful extra or mutating actions,
while retail more often appears continuation-limited under the current method.}
\label{tab:workshop_patch_families}
\end{table}

Table~\ref{tab:workshop_synthetic_results} summarizes the frozen synthetic
corpus. Heuristic search reaches perfect $24/24$ recovery with exact top-1
localization, while weaker step-order baselines degrade in predictable ways:
chronological search recovers only $7/24$ and random candidate search recovers
$15/24$.

\begin{table}[t]
\centering
\small
\begin{tabular}{lrrrr}
\toprule
Strategy & Recov. & SRR & Top-1 & MRR \\
\midrule
Heuristic & 24/24 & 1.000 & 1.000 & 1.000 \\
Reverse & 24/24 & 1.000 & 1.000 & 1.000 \\
Latest only & 24/24 & 1.000 & 1.000 & 1.000 \\
Oracle fault step & 24/24 & 1.000 & 1.000 & 1.000 \\
Random candidate & 15/24 & 0.625 & 0.583 & 0.684 \\
Chronological & 7/24 & 0.292 & 0.208 & 0.378 \\
No repair & 0/24 & 0.000 & 0.000 & 0.500 \\
\bottomrule
\end{tabular}
\caption{Workshop synthetic results on the frozen $24$-case controlled corpus.
Synthetic failures let us measure localization and recovery separately.}
\label{tab:workshop_synthetic_results}
\end{table}

These synthetic results are intentionally clean because the corruption suite is
controlled rather than ecological: each case injects a known single-step fault
into an otherwise replay-clean trajectory, primarily around tool-argument
mutations. The synthetic table is therefore a sanity-check for localization and
minimality, not a claim that natural benchmark failures are equally easy.

\subsection*{Concrete Autopsy Example}

One representative recovered retail case starts with an incorrect identity
lookup:
\begin{quote}
\small\ttfamily
find\_user\_id\_by\_email(email="aarav.santos8321@example.com")
\end{quote}
The minimal successful patch changes only the email string:
\begin{quote}
\small\ttfamily
find\_user\_id\_by\_email(email="aarav.santos8320@example.com")
\end{quote}
From that patched snapshot, strict replay with a bounded continuation proposal
reaches a successful rollout. This example is useful because it isolates the
paper's target behavior: a tiny executable patch flips the benchmark reward,
while the recovered rollout simultaneously serves as a causal autopsy. In
contrast, some airline recoveries are even smaller and succeed by deleting a
harmful mutating action altogether, showing that both continuation-style repair
and deletion-style repair appear in natural failures.

\section{Related Work}

Our work sits between benchmark evaluation for interactive agents and
post-hoc repair or deliberation methods for language models.

\paragraph{Interactive agent benchmarks.}
Recent benchmarks such as tau-style environments, WebArena, WorkArena,
AgentDojo, and SWE-bench make agent execution state, tools, and outcome
evaluation central to the task definition
\citep{tau_style,web_arena,work_arena,agent_dojo,swe_bench}. These benchmarks
are valuable because they expose realistic multi-step failures. Our goal is not
to propose another benchmark, but to turn benchmark trajectories into
debuggable execution objects by adding checkpointed intervention and replay.

\paragraph{Reasoning, reflection, and retry.}
ReAct established the interleaved reasoning-and-acting trajectory format that
many later agent systems and benchmarks inherit \citep{react}. Methods such as
Reflexion, Self-Refine, and Tree of Thoughts improve behavior through critique,
self-correction, or explicit search \citep{reflexion,self_refine,tree_of_thoughts}.
These methods usually operate in prompt space: they revise, deliberate, or
retry, but they do not center the search for the smallest executable patch that
changes the replayed benchmark outcome.

\paragraph{Intervention-supported attribution.}
The closest recent threat is REFLECT, which also uses intervention-supported
replay to strengthen attribution for silent failures in agent traces
\citep{reflect}. Our emphasis is narrower and more repair-oriented: we search
over explicit patch families, require benchmark-native reward flipping under
strict non-oracle continuation, and report patch minimality plus search cost as
first-class outputs alongside the autopsy.

\paragraph{Program repair and debugging perspective.}
The closest conceptual analogy is automated debugging or program repair, where
one searches for a small change that fixes downstream behavior. Our setting is
harder in some ways because the trajectory mixes language, tool calls, and
mutable environment state. At the same time, benchmark-native replay gives us a
clean reward signal and exact state restoration, which makes causal debugging
practical. This connects our workshop framing to classical fault localization,
delta debugging, automated program repair, and minimal counterfactual
explanations \citep{fault_localization,delta_debugging,program_repair,minimal_counterfactuals}.

\section{Artifact Summary}

The frozen workshop release now supports the following paper-facing claims and
artifacts:
\begin{itemize}
  \item real tau-native replay exists for retail and airline;
  \item the workshop corpora are frozen at $80$ natural and $24$ synthetic
  cases;
  \item saved natural failures are imported only if they remain failures under
  exact replay;
  \item strict non-oracle natural recovery is no longer only synthetic, with
  artifact-backed retail continuation recovery and airline deletion recovery;
  \item the workshop artifacts now save $27/80$ strict natural
  recoveries, $35/80$ oracle-upper-bound recoveries, and retry baselines of
  $15/80$ from scratch, $7/80$ from a localized snapshot, and $3/80$ for raw
  continuation;
  \item workshop-specific retry baselines and artifact bundling are implemented
  in the experiment harness;
  \item the workshop artifact path now includes saved oracle autopsies, paper
  tables, figure CSVs, and a curated three-case workshop panel.
\end{itemize}

The most important gaps to a full conference version are still:
\begin{itemize}
  \item hosted retry baselines at a scale beyond free-tier request caps,
  \item stronger non-oracle continuation at scale,
  \item broader baseline coverage,
  \item human evaluation of autopsy usefulness.
\end{itemize}

\section{Limitations and Next Steps}

The current system is still limited by continuation quality, patch-family
coverage, and experiment scale. Some failures are easy to localize but hard to
repair without a stronger continuation proposer. Others likely require context
editing or multi-step intervention beyond the current one-step framework. The
next milestone is to scale beyond the already-frozen $80$/$24$ workshop
release, strengthen non-oracle continuation, and evaluate whether these saved
autopsies remain useful when recovery rates drop on broader natural-failure
corpora.

\bibliographystyle{colm2026_conference}
\bibliography{references}

\end{document}
